Let $f\in\mathbb Z[x]$ be irreducible of degree $k>2$.
An integer $N\ne 0$ is called $t$-power-free if there is no prime $p$ with $p^t\mid N$.
The questions are:
(a) for $k>2$ with $k\ne 2^\ell$, does the set $\{n\in\mathbb N: f(n)\text{ is }(k-1)\text{-power-free}\}$ have positive (natural) density?
(b) are there infinitely many $n$ such that $f(n)$ is $(k-2)$-power-free?
In particular (taking $f(n)=n^4+2$ and $k=4$), are there infinitely many $n$ for which $n^4+2$ is squarefree?

Conventions/edge cases:
- We interpret $t$-power-free for negative values via $|f(n)|$.
- The density in (a) is $\lim_{x\to\infty}\frac1x\#\{1\le n\le x: f(n)\text{ is }(k-1)\text{-power-free}\}$ when the limit exists.
